\documentclass[12pt]{article}

\usepackage[margin=1in]{geometry}
\usepackage[utf8]{inputenc}
\usepackage{mathtools, amssymb, amsthm}
\usepackage{mathrsfs}
\usepackage{bbm}
\usepackage{bm}
\usepackage[colorlinks=true, linkcolor=blue, urlcolor=blue, citecolor=blue]{hyperref}
\usepackage{setspace}
\usepackage{enumitem}
\setstretch{1.25}

% required by pandoc for compact (-) lists
\providecommand{\tightlist}{%
  \setlength{\itemsep}{0pt}\setlength{\parskip}{0pt}}

\newtheorem{theorem}{Theorem}[section]
\newtheorem{lemma}[theorem]{Lemma}
\newtheorem{prop}[theorem]{Proposition}
\newtheorem{definition}[theorem]{Definition}
\newtheorem{bem}[theorem]{Bemerkung}
\newtheorem{kor}[theorem]{Korollar}
\newtheorem{deflemma}[theorem]{Definition und Lemma}

\numberwithin{equation}{section}

\renewcommand{\refname}{Referenzen}
\renewcommand{\proofname}{Beweis}

\newcommand{\ohalbe}{\frac{\omega}{2}}
\newcommand{\kk}{\cite{kk} }
\newcommand{\vol}{\operatorname{vol}}
\newcommand*{\pa}{%
    \rlap{\rotatebox{-30}{\rule[.05ex]{.4pt}{.77em}}}%
    \kern.04em%
    \rlap{\kern.36em\raisebox{0.649519052835em}{\rule{.6em}{.4pt}}}%
    \rule{.6em}{.4pt}\kern-.04em%
    \rotatebox{-30}{\rule[.05ex]{.4pt}{.77em}}}
\newcommand{\basis}{$(\omega _1, \omega _2)$ }
\newcommand{\ohnebew}{Ohne Beweis. \qed}

\title{Die Weierstraßsche \(\varphi\)-Funktion}
\author{Fabian Schuller}
\date{\today}

\begin{document}
\maketitle

\section{\texorpdfstring{Existenz der Weierstraßschen
\(\varphi\)-Funktion}{Existenz der Weierstraßschen \textbackslash varphi-Funktion}}\label{existenz-der-weierstrauxdfschen-varphi-funktion}

\begin{definition}[Gitter; \kk S. 14]

Eine Teilmenge \(\Omega \subset \mathbb C\) heißt Gitter, wenn es
reell-linear-unabhängige \(\omega _1, \omega _2 \in \mathbb C\) gibt, so
dass \begin{equation}{
\Omega = \omega _1 \mathbb Z + \omega _2 \mathbb Z.
}\end{equation}

\end{definition}

Im Folgenden sei stets \(\Omega\) ein Gitter mit Basis
\((\omega _1, \omega _2)\). Ziel ist es, eine meromorphe Funktion
\(\varphi\) zu konstruieren, die entlang des Gitters periodisch ist,
d.~h. für \(\omega \in \Omega\) und \(z \in \mathbb C\) gilt
\(\varphi(z+ \omega )=\varphi(z)\).

\begin{definition}[elliptische Funktionen]

Solche Funktionen nennt man elliptisch bezüglich \(\Omega\).

\end{definition}

\begin{theorem}[Weierstraßsche $\varphi$-Funktion; \kk S.35]\label{thm-weier}

Die Reihe \begin{equation}{
\varphi(z) := z^{-2} + \sum_{\omega \in \Omega} \left(  \frac{1}{(z-\omega )^2} - \frac{1}{\omega ^2} \right)
}\end{equation} konvergiert auf \(\mathbb C \setminus \Omega\) lokal
absolut gleichmäßig. Man nennt \(\varphi : \mathbb C \to \mathbb C\) die
Weierstraßsche \(\varphi\)-Funktion.

\end{theorem}

Für den Beweis wird die Konvergenz von Eisensteinreihen verwendet. Da es
sich um eine Summation aller Gitterpunkte handelt, erhält man mit einem
Analogon zum Riemannschen Umordnungssatz die gewünschte doppelte
Periodizität.

\begin{theorem}[Umordnungssatz; \cite{rr} S.26]\label{thm-umord}

Sei \((a_n)_{n \in \mathbb N} \subset \mathbb C\) eine Folge. Ist
\(\sum_{n=1}^\infty a_n\) absolut konvergent, so konvergiert auch jede
Umordnung dieser Reihe, d.~h. für jede Bijektion
\(k : \mathbb N \to \mathbb N\) konvergiert auch
\(\sum_{n=1}^\infty a_{k(n)}\) absolut. \ohnebew

\end{theorem}

\begin{definition}[Periodenparralelogramm, Grundmasche; \kk S. 19f]

Man definiert das Periodenparralelogramm (bezüglich
\((\omega _1, \omega _2)\)) mit Basispunkt \(u \in \mathbb C\) durch
\begin{equation}{
\pa(u, \omega _1, \omega _2) := \left\lbrace u + \alpha \omega _1 + \beta \omega _2 : \alpha, \beta  \in [0, 1) \right\rbrace
}\end{equation} und nennt
\(\pa(0, \omega_1, \omega _2) =: \pa(\omega_1, \omega _2)\) Grundmasche
von \(\Omega\).

\end{definition}

\begin{definition}[Gittervolumen; \kk S. 20]

Man definiert
\(\vol \Omega := \vol \left ( \pa(u, \omega_1, \omega _2) \right )\)
welches unabhängig von der Wahl von \(u\) und der Wahl der Basis \basis
ist. \ohnebew

\end{definition}

\begin{lemma}[\kk S.21f]\label{lem-gitter-invariant}

Man nennt \begin{equation}{
\delta := \delta (\omega _1, \omega _2) := \sup \left\{ |z-w| : z, w \in \pa(\omega _1, \omega _2)  \right\}
}\end{equation} den Durchmesser der Grundmasche von \(\Omega\). Für
\(\rho > 0\) sei \(A_\rho(\Omega)\) die Anzahl von Gitterpunkten in der
abgeschlossenen Kreisscheibe um \(0\) mit Radius \(\rho\), d.~h.
\begin{equation}{
A_\rho(\Omega) = \# \{ \omega  \in \Omega : |\omega | \le \rho \} .
}\end{equation} Dann gilt für alle \(\rho > \delta\) \begin{equation}{
\frac{\pi }{\vol \Omega} (\rho - \delta )^2 \le A_\rho(\Omega) \le \frac{\pi }{\vol \Omega} (\rho + \delta )^2 .
}\end{equation}

\begin{proof}

Betrachte \begin{equation}{
B_\rho := \{ z \in \mathbb C : |z| \le \rho \} \quad \mathrm{und} \quad M_\rho := \bigcup_{u \in \Omega : |u| \le \rho} \pa(u, \omega _1,\omega _2).
}\end{equation} Da \(\delta\) die maximale Ausbreitung einer Masche ist,
gilt \begin{equation}{
B_{\rho - \delta } \subset M_\rho \subset B_{\rho + \delta }.
}\end{equation} Geht man zum Flächenmaß dieser Mengen über erhält man
\begin{equation}{
\pi (\rho - \delta )^2 \le \vol (\Omega) A_\rho(\Omega) \le \pi (\rho + \delta )^2
}\end{equation} wegen \(\vol_2 (B_r) = \pi r^2\) und (da die
Parallelogramme an Gitterpunkten disjunkt sind) \begin{equation}{
\vol_2(M_\rho) = \sum_{u \in \Omega : |u| \le \rho} \vol_2 \left( \pa(u, \omega _1, \omega _2) \right) = \vol \Omega \cdot A_\rho(\Omega).
}\end{equation} Teilt man die Gleichung durch \(\vol(\Omega)\) erhält
man die Behauptung.

\end{proof}

\end{lemma}

\begin{deflemma}[Eisensteinreihen; \kk S.22ff]\label{thm-eisen}

Die Reihe \begin{equation}{
\sum_{0 \ne \omega  \in \Omega} |\omega|^{-\alpha}
}\end{equation} konvergiert genau dann wenn \(\alpha > 2\).
\begin{equation}{
G_k := G_k(\Omega) := \sum_{0 \ne \omega  \in \Omega} \omega^{-k}
}\end{equation} heißt \(k\)-te Eisensteinreihe (von \(\Omega\)).

\begin{proof}

Um das vorige Lemma anwenden zu können, wird die Reihe \begin{equation}{
\sum_{0 \ne \omega \in \Omega : |\omega| > \delta}  |\omega|^{-\alpha}
}\end{equation} betrachtet. Dies ist keine Einschränkung, da durch
\(|\omega| > \delta\) nur endlich viele Summanden wegfallen. Nun werden
die Summanden in die Annuli \(B_{n+1} \setminus B_n\),
\(n \in \mathbb N\) gruppiert. Eine Gruppierung kann nun mit dem Produkt
aus Anzahl von Gitterpunkten im Annulus
\(A_{n+1}(\Omega) - A_n(\Omega)\) und dem inneren oder äußeren Radius
\(n\) oder \(n+1\) abgeschätzt werden. Zwecks Abschätzung betrachten wir
jedoch Annuli der Breite \(\varepsilon > 2\delta\).

\newcommand{\sumoverannuli}{\sum_{\omega \in \Omega \cap B_{\varepsilon (n+1)} \setminus B_{\varepsilon n}}}
\newcommand{\annulicount}{ \left( A_{\varepsilon (n+1)}(\Omega) - A_{\varepsilon n}(\Omega) \right) }

\begin{equation}{
\sum_{n>\delta }^\infty \sumoverannuli  |\omega|^{-\alpha} \le \sum_{n=0}^\infty n^{-\alpha } \annulicount
}\end{equation} Aus Lemma \ref{lem-gitter-invariant} (und
\(n > \delta\)) folgt \begin{equation}{
\annulicount \le \frac{\pi }{\vol \Omega} \left(  (\varepsilon (n+1) + \delta )^2 - (\varepsilon n - \delta )^2  \right) \le c_1 n
}\end{equation} Für \(\alpha > 2\) gilt \begin{equation}{
\sum_{n>\delta }^\infty  \sumoverannuli   |\omega|^{-\alpha} \le c_1 \sum_{n=0}^\infty n^{-\alpha +1} = C < \infty
}\end{equation} mit dem Integralkriterium für Reihen. Andererseits ist
\begin{equation}{
\sum_{n>\delta }^\infty  \sumoverannuli   |\omega|^{-\alpha} \ge \sum_{n=0}^\infty {n+1}^{-\alpha }  \annulicount
}\end{equation} und analog

\begin{align}
     \annulicount &\ge \tilde c_2 \left(  (\varepsilon (n +1) - \delta )^2 - (\varepsilon n +  \delta )^2  \right) \\
&= \tilde c_2 (2n \varepsilon (\varepsilon  - \delta ) + \varepsilon^2 - 2\varepsilon \delta) \ge c_2 n
\end{align}

und man erhält Divergenz für \(\alpha \le 2\) ebenfalls mit dem
Integralkriterium. Es wurde gezeigt, dass eine Umordnung der Reihe
\begin{equation}{
\sum_{0 \ne \omega \in \Omega : |\omega| > \delta}  \omega^{-\alpha}
}\end{equation} (nicht) absolut konvergiert, und nach dem Umordnungssatz
\ref{thm-umord} also auch die Reihe selbst.

\end{proof}

\end{deflemma}

\emph{Beweis von Theorem \ref{thm-weier}.} Sei
\(K \subseteq \mathbb C \setminus \Omega\) kompakt; wähle \(R > 0\) so
dass \(K \subset B_R(0)\subset \mathbb C\). Da \(|\omega | < R+1\) für
nur endlich viele \(\omega \in \Omega\) zutrifft, kann man ohne
Einschränkung annehmen dass \(|\omega | \ge R+1\). Dann ist

\begin{align}
    \left | \frac{1}{(z-\omega )^2} - \frac{1}{\omega ^2}   \right | = \left | \frac{2z\omega - z^2 }{\omega ^2 (z-\omega )^2}   \right | = \left | \frac{2 - \frac{z}{\omega } }{(1-\frac{z}{\omega } )^2}   \right | \cdot \frac{|z|}{|\omega |^3} \le \frac{3}{(1 - \frac{R}{R+1} )^2} \cdot \frac{R}{|\omega |^3}
\end{align}

Nun folgt die Konvergenz der \(\varphi\)-Funktion aus der Konvergenz der
Eisensteinreihe \(G_3\).\qed

\begin{lemma}[\cite{kk}, S. 23]\label{lem-ungerade-eisenstein}

Es gilt \(G_k(\Omega)=0\) für alle ungeraden \(k \ge 3\).

\begin{proof}

Da für jedes \(\omega \in \Omega\) auch \(-\omega \in \Omega\), gilt
wegen des Umordnungssatzes \ref{thm-umord} \begin{equation}{
G_k(\Omega) = (-1)^k \cdot G_k(\Omega).
}\end{equation} Somit folgt die Behauptung.

\end{proof}

\end{lemma}

\begin{theorem}[Laurent-Entwicklung der Weierstraßschen $\varphi$-Funktion; \kk S.37]

Setzt man
\(\gamma := \gamma (\Omega) := \operatorname{min} \{ |\omega | : 0 \ne \omega \in \Omega \}\)
so gilt für alle \(z \in \mathbb C\) mit \(0 < |z| < \gamma\)
\begin{equation}{
\varphi (z) = z^{-2} + \sum_{n=2}^{\infty} (2n-1) G_{2n} \cdot z^{2n-2},
}\end{equation} wobei \(G_k\) die Eisenstein-Reihe aus Definition
\ref{thm-eisen} ist.

\begin{proof}

Ableiten der geometrischen Reihe ergibt \begin{equation}{
\frac{1}{(1-t)^2} = \frac{d}{dt} \left ( \frac{1}{1-t}  \right ) = \sum_{m=1}^{\infty} m t^{m-1}
}\end{equation} für \(|t|<1\), und somit gilt für \(\omega \ne 0\) und
\(|z|<\gamma\) \begin{equation}{
\frac{1}{(1-\omega)^2} - \frac{1}{\omega^2} =  \frac{1}{\omega^2} \left( \frac{1}{(1-z/\omega)^2} - 1 \right) =  \frac{1}{\omega^2} \left( \sum_{m=1}^{\infty} m \frac{z^{m-1}}{\omega^{m-1}} + 1 \right) = \sum_{m=2}^{\infty} m  \frac{z^{m-1}}{\omega^{m+1}}
}\end{equation} und daher
\begin{equation}\protect\phantomsection\label{eq:doppel}{
\varphi (z) = z^{-2} + \sum_{0 \ne \omega \in \Omega}^{} \left( \sum_{m=2}^{\infty} m  \frac{z^{m-1}}{\omega^{m+1}}   \right) , 0 < |z|<\gamma
}\end{equation} Wegen \begin{equation}{
\left| m  \frac{z^{m-1}}{\omega^{m+1}}  \right| \le \gamma m \left( \frac{z}{\gamma}  \right)^{m-1} \cdot |\omega|^{-3}
}\end{equation} kann man die Konvergenz der Eisensteinreihen
(\ref{thm-eisen}) nutzen und erhält absolute Konvergenz der Doppelreihe
in eq.~\ref{eq:doppel}. (Absolute Konvergenz in \(m\) wieder mit der
Ableitung der geometrischen Reihe.) Mit dem Umordnungssatz
\ref{thm-umord} erhält man \begin{equation}{
\varphi (z) = z^{-2} + \sum_{m=2}^{\infty} m z^{m-2} G_{m+1} , 0 < |z|<\gamma.
}\end{equation} Mit Lemma \ref{lem-ungerade-eisenstein} verschwinden die
Summanden mit \(2 \mid m\), und somit folgt die Behauptung.

\end{proof}

\end{theorem}

\section{\texorpdfstring{Differenzialgleichung für die Weierstraßsche
\(\varphi\)-Funktion}{Differenzialgleichung für die Weierstraßsche \textbackslash varphi-Funktion}}\label{differenzialgleichung-fuxfcr-die-weierstrauxdfsche-varphi-funktion}

Wir untersuchen die Weierstraßsche \(\varphi\)-Funktion mit den
Liouvilleschen Sätzen.

\begin{theorem}[Liouvillesche Sätze; \kk S.25ff]

\begin{enumerate}
\def\labelenumi{\arabic{enumi}.}
\tightlist
\item
  Ist \(f \in \mathcal K(\Omega)\) holomorph, so ist \(f\) bereits
  konstant. Hierbei bezeichnet \(\mathcal K(\Omega)\) den Körper der
  meromorphen \(\Omega\)-periodischen Funktionen.
\item
  Ist \(f \in \mathcal K(\Omega)\) und \(P\) ein
  Periodenparrallelogramm, so gilt für die Summe der Residuen
  \begin{equation}{
  \sum_{c \in P} \operatorname{res}_c f = 0
  }\end{equation}
\item
  Ist \(f \in \mathcal K(\Omega)\) nicht konstant und \(P\) ein
  Periodenparrallelogramm, so gilt für jedes \(w \in \mathbb C\)
  \begin{equation}\protect\phantomsection\label{eq:liouville-3}{
  \sum_{c \in P}^{} \operatorname{ord}_c(f-w) = 0
  }\end{equation}
\item
  Ist \(f \in \mathcal K(\Omega)\) nicht die Nullfunktion und \(P\) ein
  Periodenparrallelogramm, so gilt
  \begin{equation}\protect\phantomsection\label{eq:liouville-4}{
  \sum_{c \in P}^{} ( \operatorname{ord}_c f ) \cdot c \in \Omega
  }\end{equation}
\end{enumerate}

\ohnebew

\end{theorem}

\begin{kor}[der Laurent-Darstellung; \kk S.28]

Die Weierstraßsche \(\varphi\)-Funktion ist eine gerade Funktion, d.~h.
\(\varphi (-z)=\varphi (z)\), und für den ersten Laurent-Koeffizienten
gilt \(a_1=0\). Darüber hinaus ist \(\varphi ^\prime\) eine ungerade
(elliptische) Funktion, die genau an allen Gitterpunkten Pole dritter
Ordnung hat.

\end{kor}

\begin{lemma}[\kk S.28]\label{lem-a}

Ist \(\omega \in \Omega\) aber \(\ohalbe \notin \Omega\), dann ist
\(\ohalbe\) eine einfache Nullstelle von \(\varphi ^\prime\). Umgekehrt
gilt \(z \notin \Omega, 2z \in \Omega\) für jedes \(z \in \mathbb C\)
mit \(\varphi ^\prime(z)=0\).

\begin{proof}

Da \(\varphi ^\prime\) eine ungerade elliptische Funktion ist, gilt
\(\varphi ^\prime(z + \omega )=\varphi ^\prime(z) = - \varphi ^\prime(-z)\).
Ist \(\ohalbe\) kein Pol von \(\varphi ^\prime\) und \(\varphi\) (dies
impliziert \(\ohalbe \notin \Omega\)) gilt \begin{equation}{
\varphi ^\prime (\ohalbe) = - \varphi ^\prime(- \ohalbe + \omega) = - \varphi ^\prime(\ohalbe) = 0.
}\end{equation} Sei nun \(z_0\) eine Nullstelle von \(\varphi ^\prime\).
Wegen der Elliptizität können wir ohne Einschränkung annehmen dass
\(z_0\) in der Grundmasche \(\pa(\omega _1, \omega _2)\) liegt. Nun hat
\(\varphi ^\prime\) in der Grundmasche bereits die Nullstellen
\(\omega_1, \omega_2, \frac{\omega _1 + \omega _2}{2}\), und nach dem
dritten Liouvilleschen Satz (eq.~\ref{eq:liouville-3}) ist die Anzahl
(respektive Vielfachheit) der Nullstellen genau die Anzahl der Pole
(respektive Vielfachheit). In der Grundmasche hat \(\varphi ^\prime\)
genau einen Pol der Ordnung drei (an der Stelle \(0\)), und somit kann
es keine weiteren Nullstellen geben. Also hat \(z_0\) die gewünschte
Form.

\end{proof}

\end{lemma}

\begin{lemma}[\kk S.29]\label{lem-b}

Sei \(P\) ein Periodenparallelogramm. Für \(q \in \mathbb C\) gilt
\begin{equation}\protect\phantomsection\label{eq:ann-lemma-b}{
\forall \omega \in \Omega \quad \text{mit} \quad  \ohalbe \notin \Omega : q \ne \varphi(\ohalbe)
}\end{equation} genau dann, wenn es \(u, v \in P\) mit \(u \ne v\) aber
\(\varphi (u) = \varphi (v)=q\) gibt. In diesem Fall sind \(u\) und
\(v\) eindeutig bestimmt.

\begin{proof}

Nach dem dritten Liouvilleschen Satz (eq.~\ref{eq:liouville-3}) ist die
Anzahl der \(q\)-Stellen genau zwei. Angenommen, es gibt nur ein
\(u \in P\) mit \(\varphi (u)=q\). Dann hat die \(q\)-Stelle \(u\)
Vielfachheit zwei, und es ist \(\varphi ^\prime(u)=0\). Nach Lemma
\ref{lem-a} gilt \begin{equation}{
u \in \left\lbrace  \frac{\omega_1}{2},  \frac{\omega_2}{2}, \frac{\omega_1 + \omega_2}{2}  \right\rbrace  \mod \Omega
}\end{equation} was einen Wiederspruch zur Annahme
eq.~\ref{eq:ann-lemma-b} darstellt. Also erhalten wir \(u \ne v \in P\)
mit \(\varphi(u)=\varphi(v)=q\). Nach dem vierten Liouvilleschen Satz
(eq.~\ref{eq:liouville-4}) ist \(u+v \in \Omega\), da jeweils
\(\operatorname{ord} f=1\) gilt. Wegen dem dritten Liouvilleschen Satz
kann es auch keine weiteren \(q\)-Stellen geben.

\end{proof}

\end{lemma}

\begin{theorem}[Differenzialgleichung für die Weierstraßsche $\varphi$-Funktion; \kk S.29f]

Für alle \(z \in \mathbb C \setminus \Omega\) gilt \begin{equation}{
(\varphi ^\prime)^2(z) = 4 (\varphi(z)-e_1) (\varphi(z)-e_2) (\varphi(z)-e_3)
}\end{equation} wobei \(e_k := \varphi (\omega_k / 2)\), \(k=1,2,3\) für
eine Basis \basis von \(\Omega\) und
\(\omega_3 := \omega_1 + \omega_2\).

\begin{proof}

Nach Lemma \ref{lem-b} hat \(\varphi(z)-e_k\) eine (einzige) doppelte
Nulltelle in der Grundmasche \(P := \pa(\omega_1, \omega_2)\), nämlich
\(z_k = \omega_k / 2\). Somit besitzt auch
\(f(z) := 4 (\varphi(z)-e_1) (\varphi(z)-e_2) (\varphi(z)-e_3)\)
doppelte Nullstellen genau bei den \(z_k\). Dies gilt nach \ref{lem-a}
auch für \((\varphi ^\prime)^2(z)\), und daher ist
\(g : z \mapsto (\varphi ^\prime)^2(z) / f(z)\) eine holomorphe
(elliptische) Funktion. Der erste Satz von Liouville liefert, dass \(g\)
konstant ist. In der Nähe von \(z=0\) gilt wegen der Laurent-Entwicklung
\begin{equation}{
(\varphi ^\prime)^2(z) \approx f(z) \approx 4 (z^{-2})^3
}\end{equation} und somit erhält man \(g=1\) bei \(z \to 0\). Dies zeigt
\((\varphi ^\prime)^2 = f\) wie behauptet.

\end{proof}

\end{theorem}

\begin{thebibliography}{9}
\bibitem[KK]{kk}
Koecher, Max, und Aloys Krieg. 2007. Elliptische Funktionen und Modulformen. 2. Aufl. Masterclass. Berlin, Heidelberg: Springer Berlin Heidelberg

\bibitem[RS]{rr}
Remmert, Reinhold, und Georg Schumacher. 2002. Funktionentheorie 1. 5. Aufl. Springer-Lehrbuch. Berlin, Heidelberg: Springer Berlin Heidelberg.
\end{thebibliography}

\end{document}
